\documentclass[11pt]{article}
\usepackage[margin=1.1in]{geometry}
\usepackage{amsmath,amssymb,amsthm}
\usepackage{xcolor}
\usepackage[colorlinks=true,linkcolor=blue,citecolor=blue,urlcolor=blue]{hyperref}

\newtheorem{theorem}{Theorem}[section]
\newtheorem{proposition}[theorem]{Proposition}
\newtheorem{lemma}[theorem]{Lemma}
\newtheorem{corollary}[theorem]{Corollary}
\newtheorem{definition}[theorem]{Definition}
\theoremstyle{remark}
\newtheorem{remark}[theorem]{Remark}

\newcommand{\Z}{\mathbb{Z}}
\newcommand{\vt}[1]{v_2\!\left(#1\right)}
\newcommand{\vth}[1]{v_3\!\left(#1\right)}
\newcommand{\w}{\omega}
\newcommand{\LL}{\log_2 3}
\newcommand{\E}{\mathcal{E}_3}

\title{The 3-adic mirror of the reduced Collatz dynamics:\\ duality, dead ends, and a density bound}
\author{Ben Macindoe\thanks{\texttt{macindoebenjamin@gmail.com}. Companion to \cite{paper1}, whose coordinate system and forward laws are assumed. Developed through AI-assisted research; the methodology and verification protocol are described in Appendix~\ref{app:method}.}}
\date{July 2026}

\begin{document}
\maketitle

\begin{abstract}
The companion paper \cite{paper1} compresses the Collatz dynamics into a self-map $F$ on states $(\w,d)$, with forward per-step arithmetic governed by a $2$-adic anchor. Here we run the map backward, in three movements. First, the inverse structure: the predecessors of a state are parameterized completely by its representatives (\emph{doors}) and one free exponent $s$, and the backward depth obeys the exact law $d = 1 + \vth{s - M_3(y)}$, where the backward anchor $M_3(y)$ is an affine discrete logarithm base $2$ in the $3$-adic exponent group --- the corresponding $3$-adic mirror of the forward law $s = 2 + \vt{d - M(\w)}$, with the primes $2$ and $3$ and the coordinates $s$ and $d$ exchanged. The forward per-step theory admits a systematic $3$-adic mirror in this sense, with four structural asymmetries whose causes are identified.

Second, applications of the inverse parameterization: backward generation dies only through a single mortal door per state, dead on exactly two of the four admissible residue--parity classes; the states with no predecessors at all are precisely the depth-one states whose unique door is divisible by $3$; and \emph{steering} laws show that backward branching controls the $3$-adic data completely, freezes the $2$-adic residues, and --- through the freeze --- places the predecessor's forward anchor on an exact lattice to any prescribed precision, so that the anchor walk, unsolved in the forward direction, is controllable in reverse. Third, a density application: a deliberately impoverished sub-tree yields a self-contained proof that at least $2^{-3.6}X^{0.3}$ odd integers below $X$ reach $1$. The bound is far weaker than the linear-programming results of Krasikov--Lagarias; its contribution is a derivation in which the local branching law is exact rather than bounded by inequalities.
\end{abstract}

\section{Introduction}

Fix the odd-to-odd Collatz map $T(x) = (3x+1)/2^{\vt{3x+1}}$. The companion paper \cite{paper1} introduces reduced states $(\w, d)$ ($\w$ odd, $3\nmid\w$, $d \ge 1$), a self-map $F$ whose convergence problem is exactly equivalent to the Collatz conjecture, and shows the forward per-step arithmetic is governed by the $2$-adic anchor $M(\w) = -2\log\w/\log 9 \in \Z_2$, through the law $s = 2 + \vt{d - M(\w)}$ on the lifting classes.

This paper studies the same map in reverse, where the deterministic dynamics become an infinitely-branching tree. The conceptual center is the pair of results in Sections~\ref{sec:pred}--\ref{sec:anchor3}: the predecessors of a state admit a complete, closed-form parameterization by \emph{doors} and a single exponent (Theorem~\ref{thm:pred}), and the backward depth obeys an exact valuation law governed by a $3$-adic anchor (Theorem~\ref{thm:blaw}) that mirrors the forward law term by term. The forward theory of \cite{paper1} then admits a systematic $3$-adic mirror (Section~\ref{sec:dual}): each per-step theorem has a backward counterpart, proved here, and each counterpart differs from its original in a structural way with an identifiable cause --- the mirror is a correspondence, not an involution, and the asymmetries are part of the content.

Three consequences follow from the inverse parameterization. The \emph{dead ends} of backward generation are completely classified (Section~\ref{sec:dead}): mortality is confined to one door per state and two of the four admissible residue--parity classes, and total barrenness occurs only at depth one. The \emph{steering} laws (Section~\ref{sec:steer}) determine what infinite backward branching can control: everything $3$-adic, nothing $2$-adic directly --- and yet, through the frozen $2$-adic residues, the predecessor's \emph{forward} anchor is placed on an exact lattice, making the anchor walk controllable in reverse where it is unsolved forward. Finally, a \emph{density} application (Section~\ref{sec:density}): a single-type sub-tree, obtained through a collapse identity, yields a fully self-contained lower bound $\tilde\pi(X) \ge 2^{-3.6}X^{0.3}$ for the count of odd integers reaching $1$. This is much weaker than the computer-assisted $X^{0.84}$ of Krasikov--Lagarias \cite{kl} and is positioned accordingly; the interest is the derivation route, in which the local branching is known exactly.

\paragraph{Contributions.} (i) Complete characterization of $F^{-1}$ (Theorem~\ref{thm:pred}). (ii) The $3$-adic exponent group, the backward anchor $M_3$, and the backward valuation law (Theorem~\ref{thm:blaw}), with the exact branching ledger. (iii) The mirror correspondence: four dual theorems with proofs and identified asymmetries (Section~\ref{sec:dual}). (iv) Door mortality and the classification of predecessor-free states (Theorems~\ref{thm:mort}, \ref{thm:eden}). (v) Steering laws, including forward-anchor placement (Theorem~\ref{thm:steer}, Corollary~\ref{cor:place}). (vi) A rigorous density bound with explicit constants (Theorem~\ref{thm:density}).

\paragraph{Related work.} Predecessor trees for the $3x+1$ problem are classical: Wirsching \cite{wirsching} develops them systematically, and lower bounds on $\pi_1(x)$ run from Crandall \cite{crandall} through Krasikov's difference inequalities to the computer-assisted $x^{0.84}$ of Krasikov--Lagarias \cite{kl}; see Lagarias \cite{lagarias}. The verified range of the conjecture is $2^{71}$ \cite{barina}. A contemporary Lean-formalized conditional treatment of cycles is due to Merle \cite{merle}. Relative to this literature, the predecessor structure here is stated on reduced states rather than integers, and the local law is exact rather than an inequality; the density bound is chosen for self-containedness, not strength.

\section{Preliminaries from the companion paper}\label{sec:prelim}

For odd $x$ write $x + 1 = 2^m 3^a \w$ with $3\nmid\w$; the projection is $R(x) = (\w, m+a)$. For a state $(\w,d)$: $A = 3^d\w - 1$, $s = \vt{A}$, $x_{\mathrm{exit}} = A/2^s$, $C = A + 2^s$, $\sigma = \vt C$, $a_+ = \vth C$, and $F(\w,d) = \bigl(C/(2^\sigma 3^{a_+}),\, (\sigma - s) + a_+\bigr)$. The representatives (\emph{doors}) of $(\w,d)$ are $y_a = 2^{d-a}3^a\w - 1$, $0 \le a \le d-1$; they are pairwise distinct, the block from any representative exits at $x_{\mathrm{exit}}$, and the Collatz conjecture is equivalent to every $F$-orbit reaching $(1,1)$ \cite[Thm.~2.3]{paper1}. The forward anchor is $M(\w) = N(\w^2)$, where $9^{N(u)} = u^{-1}$ in $\Z_2$ and $\vt{9^t - 1} = 3 + \vt t$ for $t \neq 0$ (the forward isometry); $M(\w)$ is odd iff $\w \equiv \pm3 \pmod 8$ \cite[\S3]{paper1}.

\section{The predecessor structure}\label{sec:pred}

\begin{theorem}[complete characterization of $F^{-1}$]\label{thm:pred}
Let $(\Omega, D)$ be a valid state with doors $y_a$. The predecessors of $(\Omega,D)$ are exactly: for each door $y = y_a$ with $3 \nmid y$, and each $s \ge 1$ with $s$ odd if $y \equiv 1 \pmod 3$ and even if $y \equiv 2 \pmod 3$, the state
\[ (\w, d) = \bigl( N/3^{\vth N},\ \vth N \bigr), \qquad N = 2^s y + 1, \]
where $3 \mid N$ holds automatically. Distinct pairs $(a,s)$ give distinct predecessors, and there are no others. Doors with $3 \mid y$ (possible only at $a = 0$) admit no predecessors.
\end{theorem}

\begin{proof}
$F(\w,d) = (\Omega,D)$ iff $x_{\mathrm{exit}}(\w,d) = y$ for some door $y$ of $(\Omega,D)$, i.e.\ $3^d\w = 2^s y + 1$ with $s = \vt{3^d\w - 1}$. Given $(y,s)$, factoring $N = 3^d \w$ with $3\nmid\w$ forces $d = \vth N$ and $\w = N/3^d$; the divisibility $3 \mid N$ holds iff $2^s y \equiv 2 \pmod 3$, which is the stated parity condition; and $\vt{3^d\w - 1} = \vt{2^s y} = s$ exactly, since $y$ is odd.

Injectivity: suppose $(a,s)$ and $(a',s')$ produce the same predecessor state $(\w,d)$. That state has a unique exit $x_{\mathrm{exit}}(\w,d)$, which by construction equals both $y_a$ and $y_{a'}$; the doors of a state are pairwise distinct, so $a = a'$. Then $2^s y_a = 3^d\w - 1 = 2^{s'} y_a$ forces $s = s'$.

Finally, if $3 \mid y_a$ then no $s$ makes $3 \mid 2^s y_a + 1$; and $y_a \equiv 2 \pmod 3$ identically for $a \ge 1$ (see Theorem~\ref{thm:mort}), so $3 \mid y_a$ can occur only at $a = 0$.
\end{proof}

\section{The exponent group and the backward valuation law}\label{sec:anchor3}

The unit group $(\Z/3^k)^\times$ is cyclic of order $2\cdot3^{k-1}$, generated by $2$. Define the \emph{$3$-adic exponent group}
\[ \E \;:=\; \varprojlim_k \Z/(2\cdot3^{k-1})\Z \;\cong\; \Z/2 \times \Z_3, \]
so that $e \mapsto 2^e$ is an isomorphism $\E \to \Z_3^\times$. An ordinary integer $s$ embeds in $\E$ by reduction; for $e \in \E$ whose $\Z/2$ component vanishes, the $3$-adic valuation $\vth{e}$ of its $\Z_3$ component is well defined, and exponentials $2^e$ are defined by the inverse limit.

\begin{definition}\label{def:m3}
For odd $y$ with $3\nmid y$, let $M_3(y) := \log_2(-y^{-1}) \in \E$, i.e.\ the unique solution of $2^{M_3(y)} = -1/y$ in $\Z_3^\times$. Its $\Z/2$ (parity) component is determined by $y \bmod 3$ and coincides with the parity condition of Theorem~\ref{thm:pred}; its $\Z_3$ component is the backward anchor proper.
\end{definition}

\begin{proposition}[algebra of $M_3$]\label{prop:alg}
$M_3(1)$ is the distinguished exponent class representing the discrete logarithm of $-1$ to base $2$: concretely $2^{3^{k-1}} \equiv -1 \pmod{3^k}$, so $M_3(1) \equiv 3^{k-1} \pmod{2\cdot3^{k-1}}$, an element of $\E$ with odd parity component. Moreover $M_3(y_1y_2) = M_3(y_1) + M_3(y_2) - M_3(1)$: with respect to the group law of $\E$, $M_3$ is the affine logarithm $M_3(y) = M_3(1) - \log_2 y$.
\end{proposition}

\begin{proof}
$(2^{3^{k-1}})^2 = 2^{2\cdot3^{k-1}} \equiv 1 \pmod{3^k}$ and $2^{3^{k-1}} \not\equiv 1$ (its order is not a divisor of $3^{k-1}$, as $2$ generates), so $2^{3^{k-1}} \equiv -1$. The affine law is immediate from $2^{M_3(y_i)} = -y_i^{-1}$: the product of the right sides is $y_1^{-1}y_2^{-1}$, and dividing by $2^{M_3(1)} = -1$ restores the sign.
\end{proof}

\begin{lemma}[mirror isometry]\label{lem:iso3}
If $t \in \E$ has vanishing parity component and $t \neq 0$, then $\vth{2^t - 1} = 1 + \vth{t}$.
\end{lemma}

\begin{proof}
For an even integer $t$: $\vth{2^t - 1} = \vth{4^{t/2} - 1} = \vth{4 - 1} + \vth{t/2} = 1 + \vth{t}$ by lifting-the-exponent (valid since $3 \mid 4 - 1$ and $3 \nmid 4$). The statement passes to $\E$ coordinatewise: both sides depend on $t$ only through its truncations, and every parity-zero truncation is realized by an even integer.
\end{proof}

\begin{theorem}[backward valuation law]\label{thm:blaw}
For every door $y$ with $3\nmid y$ and every $s$ of the admissible parity, the element $s - M_3(y) \in \E$ has vanishing parity component, and
\[ d \;=\; \vth{2^s y + 1} \;=\; 1 + \vth{\,s - M_3(y)\,}. \]
\end{theorem}

\begin{proof}
The parity component of $M_3(y)$ equals the admissible parity of $s$ (Definition~\ref{def:m3}), so $s - M_3(y)$ is even in $\E$ and its $\Z_3$-valuation is defined. Factor
\[ 2^s y + 1 \;=\; y\,2^{M_3(y)}\bigl(2^{\,s - M_3(y)} - 1\bigr), \]
using $2^{-M_3(y)}y^{-1} = -1$; the prefactor is a $3$-adic unit. Lemma~\ref{lem:iso3} applies to the bracket.
\end{proof}

\begin{corollary}[backward ledger]\label{cor:ledger}
Along the branch family of any live door, the set of admissible $s$ with $d = j$ has exact density $2\cdot3^{-j}$ in the admissible progression: the mirror of the forward ledger $P(s = j) = 2^{-j}$.
\end{corollary}

\begin{proof}
The admissible $s$ form one arithmetic progression of step $2$; since $\gcd(2, 3^j) = 1$, among any $3^j$ consecutive admissible values, $s \bmod 3^j$ takes each residue exactly once. By Theorem~\ref{thm:blaw}, $d = j$ iff $\vth{s - M_3(y)} = j - 1$, i.e.\ iff $s \equiv M_3(y) + 3^{j-1}u \pmod{3^j}$ with $u \in \{1, 2\}$: exactly two residue classes mod $3^j$, hence exact density $2\cdot3^{-j}$.
\end{proof}

Table~\ref{tab:dual} summarizes the correspondence.

\begin{table}[h]\centering\small
\begin{tabular}{l|l}
forward \cite{paper1} & backward (this paper) \\\hline
arithmetic prime $2$ & arithmetic prime $3$ \\
$s = \vt{3^d\w - 1}$ & $d = \vth{2^s y + 1}$ \\
$M(\w) \in \Z_2$ & $M_3(y) \in \E \cong \Z/2\times\Z_3$ \\
$s = 2 + \vt{d - M(\w)}$ (lifting classes) & $d = 1 + \vth{s - M_3(y)}$ (all live doors) \\
$\vt{9^t - 1} = 3 + \vt t$ & $\vth{2^t - 1} = 1 + \vth t$ \\
ledger $2^{-j}$ & ledger $2\cdot3^{-j}$ \\
classes mod $8$ gate the law & class mod $3$ gates the parity \\
deterministic orbit & infinitely-branching tree \\
\end{tabular}
\caption{The $2\leftrightarrow3$ mirror correspondence.}\label{tab:dual}
\end{table}

\section{The mirror correspondence, with its asymmetries}\label{sec:dual}

Each per-step theorem of \cite[\S3]{paper1} has a $3$-adic mirror. We prove them and identify, in each case, the structural asymmetry. (Computational audit counts for this section appear in Appendix~\ref{app:method}, Table~\ref{tab:verify}.)

\begin{theorem}[mirror digit-determinacy]\label{thm:mdd}
(a) $M_3(y) \bmod (2\cdot 3^k)$ is determined by $y \bmod 3^{k+1}$. (b) $N = 2^s y + 1 \bmod 3^q$ is determined by $y \bmod 3^q$ and $s \bmod (2\cdot3^{q-1})$. (c) $\w = N/3^d \bmod 3^r$ is determined by $N \bmod 3^{d+r}$ and $d$. Consequently a $3$-adic window on $(y,s)$ of depth $W \ge d + r$ determines the predecessor's $\w \bmod 3^r$ exactly.
\end{theorem}

\begin{proof}
(a) $2$ has order $2\cdot3^k$ mod $3^{k+1}$, so the congruence $2^t \equiv -y^{-1} \pmod{3^{k+1}}$ determines $t$ mod $2\cdot3^k$, and $-y^{-1} \bmod 3^{k+1}$ depends only on $y \bmod 3^{k+1}$. (b) $2^s \bmod 3^q$ depends only on $s$ mod the order $2\cdot3^{q-1}$, and $y$ enters $N$ linearly. (c) Exact division by $3^d$ maps residues mod $3^{d+r}$ to residues mod $3^r$ bijectively on multiples of $3^d$. The chain is (b) $\to$ ($d$ from Theorem~\ref{thm:blaw}, visible at depth $W \geq d$) $\to$ (c).

\emph{Asymmetry.} The forward chain \cite[Thm.~3.5]{paper1} requires one genuinely cross-prime fact --- the order of $3$ mod $2^r$, needed to strip a $3$-power from a $2$-adically analyzed quantity. The backward chain requires none: $N$ is odd by construction, so the only removal ($3^d$) is same-prime relative to the residues being computed. The backward digit flow is one cross-prime step shorter.
\end{proof}

\begin{theorem}[backward increment law, and its failure set]\label{thm:minc}
Write $y_0(\kappa,K) = 2^K\kappa - 1$ for a state's $a{=}0$ door. Across a backward step through door $y$ at branch $s$, with predecessor $(\w,d)$ and $y' = y_0(\w,d)$: whenever $3 \nmid y'$, the increment $M_3(y') - M_3(y) \bmod (2\cdot3^k)$ is determined by $y \bmod 3^{d+k+1}$ and $s \bmod (2\cdot 3^{d+k})$.
\end{theorem}

\begin{proof}
$y' = 2^d\w - 1$, so $-1/y'$ is a $3$-adic function of $(\w \bmod 3^{k+1},\, 2^d \bmod 3^{k+1})$; the first input is Theorem~\ref{thm:mdd} with $r = k+1$, the second needs $d$ mod $2\cdot3^k$ (the order of $2$), and $d$ itself is exact from the window by Theorem~\ref{thm:blaw}. Then apply Theorem~\ref{thm:mdd}(a) to $y'$.

\emph{Asymmetry.} The forward increment law is total --- $F$ always produces a next state. The backward law is \emph{partial}: it is undefined exactly when $3 \mid y'$, the door-mortality set of Theorem~\ref{thm:mort}, which occupies two of the four admissible residue--parity classes. The failure is hard: no window depth resolves it, because the door is absent, not hidden.
\end{proof}

\begin{theorem}[one-step dichotomy]\label{thm:dich}
From the truncations $y \bmod 3^{K+1}$, $s \bmod (2\cdot3^K)$ alone: if $(s - M_3(y)) \not\equiv 0 \bmod 3^K$ (valuation taken in the $\Z_3$ component), the predecessor depth $d = 1 + \vth{s - M_3(y)}$ is determined exactly; otherwise the window correctly reports $d \ge K + 1$. There is no third case.
\end{theorem}

\begin{proof}
Immediate from Theorems~\ref{thm:blaw} and \ref{thm:mdd}(a): the window determines $(s - M_3(y)) \bmod 3^K$; a nonzero truncation exhibits its valuation, a zero truncation means $\vth{s - M_3(y)} \ge K$.

\emph{Asymmetry.} The forward analogue \cite[Thm.~3.6]{paper1} is a trichotomy: six of eight residue classes yield $s \in \{1,2\}$ as class constants at zero window cost, and only the lifting classes engage the window. Backward there is no free branch: Theorem~\ref{thm:blaw} is unconditional across all live doors. The cause is group-theoretic: $(\Z/3^q)^\times$ is cyclic, while $(\Z/2^q)^\times \cong \Z/2 \times \Z/2^{q-2}$; the extra $\Z/2$ factor is what furnishes the forward theory's zero-cost classes.
\end{proof}

\begin{theorem}[the two ladders]\label{thm:ladders}
\textbf{Forward.} For fixed $\w$: if $s(\w,d) = 1$ then $x_{\mathrm{exit}}(\w,d{+}1) = T(x_{\mathrm{exit}}(\w,d))$; if $s = s(\w,d) \ge 2$ then $x_{\mathrm{exit}}(\w,d{+}1) = 3\cdot2^{s-1}x_{\mathrm{exit}}(\w,d) + 1$ and $s(\w,d{+}1) = 1$.
\textbf{Backward.} For fixed door $y$, writing $N(y,s) = 3^d\w$: if $d(y,s) = 1$ then $\w(y,s{+}2) = T_3(\w(y,s))$ where $T_3(\w) = (4\w-1)/3^{\vth{4\w-1}}$, with $d(y,s{+}2) = 1 + \vth{4\w - 1}$; if $d(y,s) \ge 2$ then $\w(y,s{+}2) = 4\cdot3^{d-1}\w(y,s) - 1$ exactly, with $d(y,s{+}2) = 1$.
\end{theorem}

\begin{proof}
Forward: $A(\w,d{+}1) = 3^{d+1}\w - 1 = 3A(\w,d) + 2$. With $A = 2^s e$, $e = x_{\mathrm{exit}}$: if $s = 1$, $A' = 2(3e+1)$, so the new exit is $(3e+1)$ stripped of its $2$-power, which is $T(e)$. If $s \ge 2$, $A' = 2(3\cdot2^{s-1}e + 1)$ with odd bracket, so $s' = 1$ and $e' = 3\cdot2^{s-1}e + 1$.

Backward: $N(y,s{+}2) = 4N(y,s) - 3$. If $d \ge 2$: $N' = 3(4\cdot3^{d-1}\w - 1)$ and the bracket is $\equiv -1 \equiv 2 \pmod 3$, so $\vth{N'} = 1$ and $\w' = 4\cdot3^{d-1}\w - 1$. If $d = 1$: $N = 3\w$, so $N' = 3(4\w - 1)$, a genuine valuation collision: $\vth{N'} = 1 + \vth{4\w-1}$ and $\w' = T_3(\w)$.

\emph{Asymmetry.} Both tear-lines are gated by the respective anchors' first digits ($s = 1$ iff $\vt{d - M(\w)}$ is off-spike; $d = 1$ iff $\vth{s - M_3(y)} = 0$). The backward coefficient is $4 = 2^2$, not $3$: the parity condition confines $s$ to a lattice of index $2$, so the dual ladder's unit step is forced to be $2$. This is a consequence of already-proved structure, not a new phenomenon.
\end{proof}

\section{Dead ends}\label{sec:dead}

\begin{theorem}[door mortality]\label{thm:mort}
For $a \ge 1$, $y_a = 2^{D-a}3^a\Omega - 1 \equiv -1 \equiv 2 \pmod 3$: doors above the bottom are never dead. The sole mortal door is $a = 0$: it is dead iff $2^D\Omega \equiv 1 \pmod 3$, which holds on exactly two of the four admissible classes of $(\Omega \bmod 3,\ D \bmod 2)$.
\end{theorem}

\begin{proof}
The first claim is the displayed congruence. For $a = 0$: $3 \mid 2^D\Omega - 1$ iff $2^D\Omega \equiv 1 \pmod 3$; with $2^D \equiv (-1)^D$ and $\Omega \equiv \pm1 \pmod 3$, this selects $(\Omega \equiv 1, D \text{ even})$ and $(\Omega \equiv 2, D \text{ odd})$: two of the four classes.
\end{proof}

\begin{theorem}[predecessor-free states]\label{thm:eden}
A state has no $F$-preimage iff $D = 1$ and its unique door $2\Omega - 1$ is divisible by $3$ (equivalently $\Omega \equiv 2 \bmod 3$). Every state with $D \ge 2$ has predecessors.
\end{theorem}

\begin{proof}
By Theorem~\ref{thm:pred}, predecessors exist iff some door is live. For $D \ge 2$ the door $a = 1$ is live by Theorem~\ref{thm:mort}; for $D = 1$ the only door is $y_0 = 2\Omega - 1$, live iff $3 \nmid y_0$.
\end{proof}

Classically, one third of odd integers have no odd predecessor. In reduced coordinates the phenomenon concentrates: states of depth $\ge 2$ merely lose one door on the mortal classes, and total barrenness occurs only at depth one. The rigorous account is given in Section~\ref{sec:density}.

\section{Steering}\label{sec:steer}

Fix a live door $y$ and sweep the admissible $s$.

\begin{theorem}[steering laws]\label{thm:steer}
(i) \emph{Depth: total control.} For every $d^* \ge 1$, the branches with $d = d^*$ form a subset of exact density $2\cdot3^{-d^*}$ of the admissible $s$; in particular every depth is realized by infinitely many branches.
(ii) \emph{$2$-adic residues: frozen.} For every $k$ and every admissible $s \ge k$, the predecessor satisfies $\w \equiv 3^{-d} \pmod{2^{k}}$. Direct $2$-adic steering is impossible beyond the finitely many branches with $s < k$.
(iii) \emph{Forward-anchor placement.} For every admissible $s \ge 3$, the predecessor's forward anchor satisfies
\[ M(\w) \;\equiv\; d \pmod{2^{\,s-2}}, \]
where $d = d(y,s)$. (Equivalently $M(\w) \equiv -d\,M(3)$, since $M(3) = N(9) = -1$.)
\end{theorem}

\begin{proof}
(i) is Corollary~\ref{cor:ledger}. (ii): $N = 2^s y + 1 \equiv 1 \pmod{2^s}$, so $\w = N\cdot3^{-d} \equiv 3^{-d} \pmod{2^{\min(s,k)}}$. (iii): from $\w\,3^d = 1 + 2^s y$,
\[ \w^2 9^{\,d} = (1 + 2^s y)^2 = 1 + 2^{s+1}y\,(1 + 2^{s-1}y), \]
with odd second factor, so $\vt{\w^29^{\,d} - 1} = s + 1 \ge 4$: the argument lies in $1 + 8\Z_2$, where $N$ is defined, and $N$ is a homomorphism there. Thus $N(\w^2 9^{\,d}) = M(\w) + d\,N(9) = M(\w) - d$, since $9^{N(9)} = 9^{-1}$ gives $N(9) = -1$. The forward isometry gives $\vt{N(u)} = \vt{u - 1} - 3$, so $\vt{M(\w) - d} = s - 2$, which is the display (with equality of valuation, so the modulus is sharp).
\end{proof}

\begin{remark}[relation to the forward valuation law]\label{rem:samelaw}
Unwound, Theorem~\ref{thm:steer}(iii) is not a new identity: since $\w\,3^d = 1 + 2^s y$ is precisely the exit equation of the forward step, the sharp congruence $\vt{M(\w) - d} = s - 2$ is the forward valuation law $s = 2 + \vt{d - M(\w)}$ of \cite{paper1}, encountered from the other side. What is new is the reading: forward, the state $(\w, d)$ is given and the law \emph{reveals} $s$; backward, the pair $(y, s)$ is chosen freely and the same law \emph{places} $d$ --- and with it, by this congruence, the predecessor's anchor residue. The two papers' central laws are one identity viewed from opposite ends of the step.
\end{remark}

\begin{corollary}[anchor placement to arbitrary precision]\label{cor:place}
For the fixed live door $y$: for every $k$ and every target residue $\rho \bmod 2^k$ there exist infinitely many admissible $s$ whose predecessor satisfies $M(\w) \equiv \rho \pmod{2^k}$: choose any $d^* \ge 1$ with $d^* \equiv \rho \pmod{2^k}$, apply Theorem~\ref{thm:steer}(i) to obtain infinitely many admissible $s \ge k + 2$ with $d = d^*$, and Theorem~\ref{thm:steer}(iii) to each.
\end{corollary}

\begin{remark}
Forward, the anchor increment along an orbit is the central unresolved object of \cite{paper1}. Backward, by Corollary~\ref{cor:place}, the anchor is placeable at will. The same walk changes epistemic status with direction of travel; whether backward controllability yields forward information is, in our judgment, the sharpest question this paper raises.
\end{remark}

\section{A density bound}\label{sec:density}

\begin{definition}[door tree]\label{def:doortree}
Root $y = 1$. Children of a node $y$ (odd, $3\nmid y$): for each admissible $s$ (parity per $y \bmod 3$), the integer $y' = (2^{s+1}y - 1)/3$, kept when $3\nmid y'$ and $y' > 1$.
\end{definition}

\begin{lemma}[collapse identity]\label{lem:collapse}
For every admissible $s$, $3 \mid 2^{s+1}y - 1$. Let $(\w', d)$ be the predecessor of $y$'s state through $(y,s)$. If $d = 1$, its unique door is $2\w' - 1 = (2^{s+1}y-1)/3$; if $d \ge 2$, its door at $a = d-1$ is $2\cdot3^{d-1}\w' - 1 = (2^{s+1}y-1)/3$: \emph{the same value}. The door map $y \mapsto (2^{s+1}y-1)/3$ is independent of the predecessor's depth.
\end{lemma}

\begin{proof}
Admissibility gives $2^s y \equiv 2 \pmod 3$, hence $2^{s+1}y \equiv 1$ and $3 \mid 2^{s+1}y - 1$. For $d = 1$: $\w' = (2^sy+1)/3$ and $2\w' - 1 = (2^{s+1}y + 2 - 3)/3$. For $d \ge 2$: $\w' = (2^sy+1)/3^d$ and $2\cdot3^{d-1}\w' - 1 = 2(2^sy+1)/3 - 1$, the same value.
\end{proof}

\begin{lemma}[triple law]\label{lem:triple}
For any three consecutive admissible $s$, the values $2^s y + 1$ are congruent to $0, 3, 6 \pmod 9$ in some order. Consequently the three candidate $y'$ are $\equiv 0, 1, 2 \pmod 3$ in some order: exactly two children per triple are kept (as live doors), one of each depth type.
\end{lemma}

\begin{proof}
Consecutive admissible $s$ differ by $2$, and $2^{s+2}y + 1 - (2^s y + 1) = 3\cdot2^s y \equiv 6 \pmod 9$ (using $2^s y \equiv 2 \bmod 3$). So the three values are $v, v+6, v+12 \equiv v, v+6, v+3 \pmod 9$: distinct, and all divisible by $3$; hence $\{0,3,6\}$. The value $\equiv 3$ gives $\vth N = 1$ and $\w' \equiv 1 \pmod 3$, a live depth-$1$ child; the value $\equiv 0$ gives $d \ge 2$, whose door is $\equiv 2 \pmod 3$ (Theorem~\ref{thm:mort}), live; the value $\equiv 6$ gives a dead door.
\end{proof}

\begin{lemma}[unique parent, distinct values]\label{lem:distinct}
Every kept child satisfies $y' > y$, except that the root's children satisfy $y' \ge 5$. Each node value determines its generating edge: $y'$ determines its state (if $y' \equiv 1 \bmod 3$ the state is $((y'+1)/2, 1)$; if $y' \equiv 2 \bmod 3$, then for generated nodes $y'$ is the terminal door of its state, so $d = 1 + \vth{(y'+1)/2}$ and $\w = (y'+1)/(2\cdot3^{d-1})$), the state determines the parent door $y = x_{\mathrm{exit}}$ of that state, and then $s$ is determined by $2^{s+1}y = 3y' + 1$. Consequently all node values in the door tree are pairwise distinct integers.
\end{lemma}

\begin{proof}
Growth: for $y > 1$ and $s \ge 1$, $y' = (2^{s+1}y - 1)/3 \ge (4y-1)/3 > y$; at the root, the smallest kept child is $y' = 5$ (at $s = 3$; $s = 1$ gives the excluded self-loop $y' = 1$). Unique generation: as stated --- the state recovery is \cite[\S2]{paper1} applied to $y' + 1$; the parent door is the exit of the child's state by Theorem~\ref{thm:pred}; and $s$ is pinned by the displayed equation. Distinctness follows by induction along the (finite, by growth) ancestry of any node: two nodes of equal value would have equal generating edges, hence equal parents of equal value, and so on back to the root, forcing them to be the same node. Growth also excludes any directed cycle: values strictly increase along edges away from the root, and the only cycle of the unrestricted door relation at the root is the excluded $s{=}1$ self-loop.
\end{proof}

\begin{lemma}[mass, non-root]\label{lem:mass}
Let $y > 1$ be a node and let $I(y)$ be the kept children of its first two windows of three consecutive admissible $s$ each (i.e.\ $s \le s_0 + 10$, where $s_0 \le 2$ is the least admissible exponent). Then every $z \in I(y)$ satisfies $z < 2^{11.415}\,y$, and at $c = 0.3$,
\[ \sum_{z \in I(y)} (z/y)^{-c} \;\ge\; 2^{-3.415c} + 2^{-5.415c} + 2^{-9.415c} + 2^{-11.415c} \;=\; 1.0502\ldots \;>\; 1. \]
\end{lemma}

\begin{proof}
Each kept child at branch $s$ satisfies $z < 2^{s+1}y/3 = 2^{\,s+1-\LL}y$, so its log-increment is below $s - 0.585$; and $z > 2^{s-1}y \ge y$ gives the crude lower bound used only for ordering. By Lemma~\ref{lem:triple}, each window of three consecutive admissible $s$ (spanning six integers) contains two kept children. The adversarial placement (kept children at the two largest slots of each window, worst offset $s_0 = 2$) yields increments at most $3.415,\,5.415$ (first window) and $9.415,\,11.415$ (second), giving the displayed sum; any true placement dominates it termwise. The largest possible child in these windows has $s \le 12$, hence $z < 2^{13 - \LL}y = 2^{11.415}y$.
\end{proof}

\begin{lemma}[mass, root]\label{lem:rootmass}
The root's kept children at $s = 3, 7, 9$ are $5, 85, 341$, all below $2^{12}$, and at $c = 0.3$: $5^{-0.3} + 85^{-0.3} + 341^{-0.3} = 1.0546\ldots > 1$.
\end{lemma}

\begin{proof}
Direct computation ($s = 5$ and $s = 11$ give the dead doors $21$ and $1365$).
\end{proof}

\begin{lemma}[renewal induction]\label{lem:renewal}
Let $c = 0.3$ and $A = 2^{-12c}$. For every node $y$ of the door tree and every $X \ge 1$, the number $f(y, X)$ of descendants of $y$ (inclusive) with value $\le X$ satisfies $f(y,X) \ge A\,(X/y)^c$ whenever $X \ge y$.
\end{lemma}

\begin{proof}
Strong induction on $\lceil X/y \rceil$. If $X < 2^{12} y$: $f \ge 1 \ge A(X/y)^c$ since $A\,2^{12c} = 1$. If $X \ge 2^{12}y$: every $z \in I(y)$ (Lemma~\ref{lem:mass}; at the root, the three children of Lemma~\ref{lem:rootmass}) satisfies $z < 2^{12} y \le X$. Moreover every child satisfies $z \ge \tfrac{19}{15}y$ (non-root: $z \ge (4y-1)/3$ with $y \ge 5$; root: $z \ge 5$), so $X/z \le \tfrac{15}{19}(X/y)$, and since $X/y \ge 2^{12}$ the drop exceeds $1$: $\lceil X/z \rceil < \lceil X/y \rceil$, and the inductive hypothesis applies to each child. The descendant sets of distinct children are disjoint (Lemma~\ref{lem:distinct}). Hence
\[ f(y,X) \;\ge\; \sum_{z \in I(y)} f(z, X) \;\ge\; A X^c \sum_{z\in I(y)} z^{-c} \;\ge\; A\,(X/y)^c \cdot 1.05 \;\ge\; A\,(X/y)^c. \qedhere \]
\end{proof}

\begin{theorem}[density bound]\label{thm:density}
Let $\tilde\pi(X) = \#\{\text{odd } y \le X : \text{the } T\text{-orbit of } y \text{ reaches } 1\}$. For all $X \ge 1$,
\[ \tilde\pi(X) \;\ge\; 2^{-3.6}\, X^{0.3}. \]
The same argument yields any exponent $c$ with $(2^{-3.415c} + 2^{-5.415c})/(1 - 2^{-6c}) > 1$ (critical value $c^* \approx 0.3304$), with correspondingly adjusted constant.
\end{theorem}

\begin{proof}
Nodes of the door tree are doors of states backward-reachable from $(1,1)$; by the equivalence theorem of \cite{paper1}, each node's $T$-orbit reaches $1$. Nodes are pairwise distinct odd integers (Lemma~\ref{lem:distinct}). Apply Lemma~\ref{lem:renewal} at the root: $\tilde\pi(X) \ge f(1, X) \ge 2^{-3.6}X^{0.3}$. For general $c < c^*$: take sufficiently many non-root windows that the finite partial mass exceeds $1$ at that $c$, and sufficiently many root children likewise; the induction of Lemma~\ref{lem:renewal} then runs verbatim with a correspondingly larger scale threshold and smaller constant $A$.
\end{proof}

\begin{remark}[position]
The exponent sits between Crandall \cite{crandall} and Krasikov's 1989 value $0.43$; Krasikov--Lagarias reach $0.84175$ by linear programming over difference-inequality systems mod $3^{11}$ \cite{kl}. The comparison is deliberate: the core above uses one door per state, two children per triple, and adversarial anchor phases; empirically the same core grows at exponent $\approx 0.45$, and the full backward tree at $\approx 0.98$. The discarded resources have close analogues in the residue refinements and optimization stages of the Krasikov--Lagarias program, suggesting a possible route from the present branch inventory to a richer finite-type counting system; constructing that translation, and determining whether an exact local law improves on inequalities as input to it, are open problems, claimed in neither direction.
\end{remark}

\section{Discussion}

The mirror correspondence completes a symmetry the companion paper only gestured at: forward dynamics consume $2$-adic digits under exact laws with no choice; backward dynamics offer infinite choice that moves only $3$-adic data, with the $2$-adic side frozen --- yet the freeze routes complete control of the forward anchor through the choice of depth (Corollary~\ref{cor:place}). The conjecture, in these coordinates, is the assertion that a deterministic $2$-adic flow and a controllable $3$-adic tree are the same object. The asymmetries of Section~\ref{sec:dual} say the two sides are not interchangeable, and in several respects the backward theory is simpler: one fewer cross-prime step, no class trichotomy, an unconditional valuation law. The concrete open items, in increasing ambition: enrich the density core with doors, residue types, and anchor phases, and optimize; make the supercritical renewal heuristic of Section~\ref{sec:dead} rigorous (a mortality-corrected branching heuristic for the full tree, with the measured stationary depth law as its one non-rigorous input, gives a renewal mass exceeding $1$ at every exponent below $1$, consistent with exact enumeration of the tree to $\w \le 2^{19}$; we do not use this heuristic anywhere in the paper, and flag it as such); and determine whether backward anchor controllability implies anything about the forward anchor walk.

\appendix
\section{Methodology, responsibility, and verification}\label{app:method}
This paper continues the protocol of \cite[App.~A]{paper1}: the author directs and takes responsibility; the mathematical development is AI-assisted; every result was tested with independently written code \emph{before} proof finalization, and the proofs above are intended to stand without the computations. One section (the four mirror theorems of Section~\ref{sec:dual}) was first developed in a delegated session of a different model under a written brief with branch discipline, then independently re-verified and proof-audited before inclusion. Two corrections occurred during this paper's work and are part of the record: an initial renewal-equation heuristic that ignored representative multiplicity, and a conjectured one-step window-steering property refuted by its own first test (the refutation became Theorem~\ref{thm:steer}(ii)). The complete research record, including all verification scripts, is public at \url{https://github.com/macindoe/collatz}.

\begin{table}[h]\centering\small
\begin{tabular}{l|l}
result & computational audit (all zero failures) \\\hline
Theorem~\ref{thm:pred} & exact match vs.\ brute-force forward scans, $7$ targets \\
Theorem~\ref{thm:blaw} & $4{,}265$ random $(y,s)$ at anchor depth $3^8$ \\
Corollary~\ref{cor:ledger} & measured $0.6664/0.2230/0.0736/\ldots$ vs.\ $2/3, 2/9, 2/27,\ldots$ \\
Theorem~\ref{thm:mdd} & $3{,}000$ checks per fact; $\approx2{,}670$ per window depth \\
Theorem~\ref{thm:minc} & $6{,}000$ trials; failure set matches mortality exactly \\
Theorem~\ref{thm:dich} & $\approx13{,}200$ trials per $K \in \{2,4,6,8\}$ \\
Theorem~\ref{thm:ladders} & $30{,}000$ states (forward), $19{,}992$ (backward) \\
Theorems~\ref{thm:mort}, \ref{thm:eden} & $20{,}000$ states all doors; $600$ states vs.\ forward image \\
Theorem~\ref{thm:steer}(iii) & $2{,}025$ checks at depth $20$; bound $s-2$ attained exactly \\
Lemmas~\ref{lem:collapse}, \ref{lem:triple} & $13{,}408$ and $20{,}000$ checks \\
door tree & built to $2^{18}$; all sampled nodes reach $1$; exponent $\approx0.45$ \\
\end{tabular}
\caption{Computational audits (methodology information, not proof).}\label{tab:verify}
\end{table}

\begin{thebibliography}{9}
\bibitem{paper1} B.~Macindoe, \emph{Reduced coordinates for the Collatz map: exact per-step laws, anchor dynamics, and the limits of counting arguments for cycles}, preprint, Zenodo DOI 10.5281/zenodo.21273548 (2026).
\bibitem{crandall} R.~E.~Crandall, \emph{On the $3x+1$ problem}, Math.\ Comp.\ 32 (1978), 1281--1292.
\bibitem{kl} I.~Krasikov, J.~C.~Lagarias, \emph{Bounds for the $3x+1$ problem using difference inequalities}, Acta Arith.\ 109 (2003), 237--258; arXiv:math/0205002.
\bibitem{wirsching} G.~J.~Wirsching, \emph{The Dynamical System Generated by the $3n+1$ Function}, Lecture Notes in Math.\ 1681, Springer (1998).
\bibitem{lagarias} J.~C.~Lagarias, \emph{The 3x+1 problem: an annotated bibliography}, I--II, arXiv math/0309224, math/0608208.
\bibitem{barina} D.~Barina, \emph{Improved verification limit for the convergence of the Collatz conjecture}, J.\ Supercomputing (2025).
\bibitem{merle} E.~Merle, \emph{On the non-existence of non-trivial Collatz cycles: a conditional formal proof in Lean 4}, preprint, Zenodo DOI 10.5281/zenodo.19790406 (2026).
\end{thebibliography}

\end{document}
